\chapter{Epididymitis}

\section*{Definition and Overview}
Epididymitis is inflammation of the epididymis, the coiled tube at the back of the testicle that stores and carries sperm. It is usually caused by infection, which may spread from the urinary tract or, in younger sexually active men, be a sexually transmitted infection. The condition causes pain and swelling in the scrotum, sometimes with fever, and it requires medical assessment to distinguish it from more urgent conditions such as testicular torsion. Treatment is with antibiotics, pain relief, and rest, and most men recover fully.

\section*{Epidemiology}
Epididymitis is one of the most common causes of scrotal pain in adult men. In sexually active men under 35, the most common cause is a sexually transmitted infection such as chlamydia or gonorrhoea, while in older men it is usually caused by urinary tract bacteria. The condition can occur at any age and is more common with urinary problems such as an enlarged prostate, catheter use, or recent urinary procedures.

\section*{Aetiology and Pathophysiology}
Epididymitis is caused by infection or inflammation of the epididymis.
\begin{itemize}
    \item \textbf{STIs:} chlamydia and gonorrhoea, usually in sexually active men under 35
    \item \textbf{Urinary bacteria:} Escherichia coli and other bacteria from the urinary tract, more common in older men
    \item \textbf{Spread:} bacteria travel backwards up the tubes from the prostate and urinary tract
    \item \textbf{Other causes:} reflux of urine, urine stasis, trauma, and some medicines
    \item \textbf{Risk factors:} urinary obstruction, catheter use, recent urological procedures, and unprotected sex
    \item \textbf{Mechanism:} inflammation causes swelling of the epididymis, and the infection can spread to involve the testicle itself, called epididymo-orchitis
    \item \textbf{Chronic epididymitis:} persistent or recurrent inflammation without infection
\end{itemize}

\section*{Clinical Features}
Symptoms usually develop gradually over a few days.
\begin{itemize}
    \item Pain and tenderness in the scrotum, often on one side
    \item Swelling and redness of the scrotum
    \item A feeling of heaviness in the testicle
    \item Fever and chills in some cases
    \item Pain on urination, urinary frequency, or discharge
    \item Pain relieved by lifting the testicle
    \item In chronic epididymitis: persistent, milder discomfort
\end{itemize}

\section*{Differential Diagnosis}
Scrotal pain has causes that need urgent distinction.
\begin{itemize}
    \item Testicular torsion, a surgical emergency in which the blood supply to the testicle is cut off, causing sudden severe pain
    \item Orchitis, inflammation of the testicle itself
    \item Testicular cancer, which can cause painless swelling or heaviness
    \item Inguinal hernia
    \item Kidney stones causing referred pain
    \item Hydrocele and varicocele
\end{itemize}

\section*{When to Seek Medical Care}
Any man with scrotal pain or swelling should be assessed urgently by a doctor, because testicular torsion must be excluded without delay.

\textbf{Call 000 (or local emergency services) for:}
\begin{itemize}
    \item Sudden severe scrotal pain
    \item Scrotal pain with nausea and vomiting
    \item A swollen, tender scrotum in a boy, which may be torsion
    \item Fever with severe pain
    \item Collapse or signs of sepsis
\end{itemize}

\section*{Diagnosis and Investigations}
Diagnosis is made by examination and testing.
\begin{itemize}
    \item \textbf{Examination:} assessment of the scrotum, and the position and tenderness of the testicle
    \item \textbf{Urine tests:} urinalysis and culture
    \item \textbf{STI testing:} urine or swab tests for chlamydia and gonorrhoea in sexually active men
    \item \textbf{Ultrasound:} duplex ultrasound assesses blood flow and distinguishes epididymitis from torsion
    \item \textbf{Blood tests:} inflammatory markers and full blood count
\end{itemize}

\section*{Management}
Treatment is with antibiotics and supportive measures.
\begin{itemize}
    \item \textbf{Antibiotics:} directed at the likely organism, adjusted to test results
    \item \textbf{Partner treatment:} for STI-related epididymitis, sexual partners are treated and tested
    \item \textbf{Pain relief:} anti-inflammatory medicines
    \item \textbf{Scrotal support:} supportive underwear and elevation of the scrotum
    \item \textbf{Rest:} avoiding heavy activity until the pain settles
    \item \textbf{Cold packs:} applied to the scrotum for swelling
    \item \textbf{Follow-up:} to confirm resolution and check for complications
    \item \textbf{Chronic epididymitis:} pain management, and investigation of underlying causes
\end{itemize}

\section*{Complications}
\begin{itemize}
    \item Spread of infection to the testicle, causing epididymo-orchitis
    \item Abscess formation
    \item Chronic pain and persistent discomfort
    \item Scarring of the epididymis, which can affect fertility
    \item Sepsis in severe cases
    \item Recurrence
\end{itemize}

\section*{Prognosis and Outcomes}
With prompt antibiotic treatment, most men with epididymitis improve within a few days and recover fully within a few weeks. STI-related cases respond well when partners are treated to prevent reinfection. Chronic epididymitis is more difficult to treat, but pain can usually be managed. Fertility is rarely affected, but complications should be treated promptly.

\section*{Prevention and Self-Care}
\begin{itemize}
    \item Practise safe sex and get tested for sexually transmitted infections
    \item Treat urinary infections promptly
    \item Seek early care for any scrotal pain
    \item Wear supportive underwear during the illness
    \item Rest and avoid strenuous activity
    \item Complete the full course of antibiotics
    \item Attend follow-up, especially if symptoms persist
\end{itemize}

\section*{Sources and Further Reading}
The following reputable sources provide further information for healthcare students and patients seeking reliable, current advice about epididymitis.
\begin{itemize}
    \item Healthdirect Australia. \textit{Epididymitis.} https://www.healthdirect.gov.au/epididymitis
    \item Urology Care Foundation. \textit{Epididymitis patient information.} https://www.urologyhealth.org/
    \item NHS. \textit{Epididymitis and orchitis.} https://www.nhs.uk/conditions/epididymitis-orchitis/
    \item Mayo Clinic. \textit{Epididymitis: symptoms and causes.} https://www.mayoclinic.org/diseases-conditions/epididymitis/
    \item Healthy Male Australia. \textit{Testicular and scrotal conditions.} https://www.healthymale.org.au/
    \item MSD Manual. \textit{Epididymitis: professional overview.} https://www.msdmanuals.com/
\end{itemize}
